% Part: model-theory
% Chapter: models-of-arithmetic
% Section: non-standard-models

\documentclass[../../../include/open-logic-section]{subfiles}

\begin{document}

\olfileid{mod}{mar}{mdq}
\section{$\Th{Q}$ నమూనాలు}

\begin{explain}
$\Struct{N}$ సత్యం చేసే \tetoken{వాక్యాలనే}{!!{sentence}s} సత్యం
చేసే అప్రామాణిక \tetoken{నిర్మాణాలు}{!!{structure}s} ఉన్నాయని మనకు
తెలుసు; అంటే, $\Th{TA}$కు అప్రామాణిక నమూనా ఉంది. $\Sat{N}{\Th{Q}}$
కాబట్టి $\Th{TA}$ యొక్క ప్రతి నమూనా $\Th{Q}$కు కూడా నమూనా.
$\Th{Q}$, $\Th{TA}$ కంటే చాలా బలహీనమైనది; ఉదాహరణకు
$\Th{Q}\Proves/\lforall[x][\lforall[y][\eq[(x+y)][(y+x)]]]$.
బలహీన సిద్ధాంతాలను సంతృప్తిపరచడం సులభం: వాటికి ఎక్కువ నమూనాలు
ఉంటాయి. ఉదాహరణకు
$\lforall[x][\lforall[y][\eq[(x+y)][(y+x)]]]$ను అసత్యం చేసే
$\Th{Q}$ నమూనాలు ఉన్నాయి; కానీ అవి $\Th{TA}$కు, లేదా
$\Th{PA}$కు కూడా, నమూనాలు కాలేవు. $\Th{Q}$ నమూనాలు సాపేక్షంగా
సరళమైనవి కూడా: వాటిని స్పష్టంగా నిర్దేశించవచ్చు.
\end{explain}

\begin{ex}
\ollabel{ex:model-K-of-Q}
$\Domain{K}=\Nat\cup\{a\}$ గల \tetoken{నిర్మాణం}{!!{structure}}~
$\Struct{K}$ను, ఈ అర్థనిర్దేశాలతో పరిశీలించండి:
\begin{align*}
  \Assign{\Obj{0}}{K} & = 0\\
  \Assign{\prime}{K}(x) & =
  \begin{cases}
    x+1 & \text{$x\in\Nat$ అయితే}\\
    a & \text{$x=a$ అయితే}
  \end{cases}\\
  \Assign{+}{K}(x, y) & =
  \begin{cases}
    x+y & \text{$x,y\in\Nat$ అయితే}\\
    a & \text{మిగతా సందర్భాల్లో}
  \end{cases}\\
  \Assign{\times}{K}(x, y) & =
  \begin{cases}
    xy & \text{$x,y\in\Nat$ అయితే}\\
    0 & \text{$x=0$ లేదా $y=0$ అయితే}\\
    a & \text{మిగతా సందర్భాల్లో}\\
  \end{cases}\\
  \Assign{<}{K} & =
  \Setabs{\tuple{x,y}}{x, y \in \Nat \text{ మరియు } x<y} \cup
  \Setabs{\tuple{x,a}}{x \in \Domain{K}}
\end{align*}
$\Sat{K}{\Th{Q}}$ అని చూపాలంటే, $\Th{Q}$ స్వీకృతాలన్నీ
$\Struct{K}$లో సత్యమని నిర్ధారించాలి. సౌలభ్యం కోసం
$\Assign{\prime}{K}(x)$ను ($\Struct{K}$లో $x$ యొక్క
“ఉత్తరాధికారి”) $x^\nssucc$గా, $\Assign{+}{K}(x,y)$ను
($\Struct{K}$లో $x$, $y$ల “మొత్తం”) $x\nsplus y$గా,
$\Assign{\times}{K}(x,y)$ను ($\Struct{K}$లో $x$, $y$ల “లబ్ధం”)
$x\nstimes y$గా, $\tuple{x,y}\in\Assign{<}{K}$ను
$x\nsless y$గా రాద్దాం. ఈ సంక్షిప్తీకరణలతో $\Struct{K}$లోని
క్రియలను మరింత స్పష్టంగా ఇలా ఇవ్వవచ్చు:
\[
\begin{array}{c|c}
  x & x^\nssucc \\
  \hline
  n & n+1 \\
  a & a
\end{array}
\qquad
\begin{array}{c|ccc}
  x \nsplus y & 0 & m & a \\
  \hline
  0 & 0 & m & a \\
  n & n & n+m & a \\
  a & a & a & a \\
\end{array}
\qquad
\begin{array}{c|ccc}
  x \nstimes y & 0 & m & a \\
  \hline
  0 & 0 & 0 & 0 \\
  n & 0 & nm & a \\
  a & 0 & a & a \\
\end{array}
\]
$n$, $m\in\Nat$లకు $n\nsless m$ అవడం సరిగ్గా $n<m$ అయినప్పుడే;
ప్రతి $x\in\Domain{K}$కూ $x\nsless a$.

$\nssucc$ \tetoken{ఒకటి-ఒకటి}{!!{injective}} కాబట్టి
$\Sat{K}{\lforall[x][\lforall[y][(\eq[x'][y']\lif\eq[x][y])]]}$.
$\Struct{K}$లో $0$ ఏ $\nssucc$-ఉత్తరాధికారీ కాదు కాబట్టి
$\Sat{K}{\lforall[x][\eq/[\Obj 0][x']]}$. ప్రతి $n>0$కు
$n=(n-1)^\nssucc$, అలాగే $a=a^\nssucc$ కాబట్టి
$\Sat{K}{\lforall[x][(\eq[x][\Obj 0]\lor
\lexists[y][\eq[x][y']])]}$.

$n\nsplus0=n+0=n$, $\nsplus$ నిర్వచనం వల్ల $a\nsplus0=a$ కాబట్టి
$\Sat{K}{\lforall[x][\eq[(x+\Obj 0)][x]]}$.
$\Sat{K}{\lforall[x][\lforall[y][\eq[(x+y')][(x+y)']]]}$ కొంచెం
క్లిష్టమైనది. $n$, $m$ రెండూ ప్రామాణికమైతే
\begin{align*}
(n \nsplus m^\nssucc) & = (n+(m+1)) = (n+m)+1 = (n \nsplus m)^\nssucc
\intertext{ఎందుకంటే ప్రామాణిక సంఖ్యలపై $\nsplus$, $^\nssucc$లు
  $+$, $\prime$లతో ఏకీభవిస్తాయి. ఇప్పుడు
  $x\in\Domain{K}$ అనుకుందాం. అప్పుడు}
(x \nsplus a^\nssucc) & = (x \nsplus a) = a = a^\nssucc
  = (x \nsplus a)^\nssucc
\intertext{మిగిలిన సందర్భంలో $x=a$, $y\in\Domain{K}$. ఇక్కడ
  $y=n$ ప్రామాణికమా, లేక $y=a$నా అనేదాన్ని కూడా వేరు చేయాలి:}
(a \nsplus n^\nssucc) & = (a \nsplus (n+1)) = a = a^\nssucc
  = (a \nsplus n)^\nssucc\\
(a \nsplus a^\nssucc) & = (a \nsplus a) = a = a^\nssucc
  = (a \nsplus a)^\nssucc
\end{align*}
\sourcecorrection{OLTEMODARI-008}{$\Struct{K}$ వ్యక్తి క్షేత్రంలో
$\Nat$కు అదనంగా ఉన్నది $a$ మాత్రమే; అయినా ఆంగ్ల మూలం చివరి
సందర్భ విభజనలో $y=b$ అంటుంది. దాని వెంటనే ఉన్న సరైన $a$-సూత్రాలకు
అనుగుణంగా $y=a$గా సరిచేశాం.}
ఇది అవసరమైనదానికంటే కొంచెం వివరంగానే ఉంది. ఉదాహరణకు,
$z$ ఏదైనా $a\nsplus z=a$ కాబట్టి
$a\nsplus a^\nssucc=a$ అని వెంటనే చెప్పవచ్చు. మిగిలిన స్వీకృతాలను
కూడా ఇదే విధంగా నిర్ధారించవచ్చు.

ఈ విధంగా $\Struct{K}$, $\Th{Q}$కు నమూనా. దాని “సంకలనం”
$\nsplus$ వినిమయ ధర్మం గలది కూడా. అయితే $\Struct{N}$లో సత్యమై
$\Struct{K}$లో అసత్యమైన, అలాగే విపర్యయంగా ఉండే ఇతర వాక్యాలు ఉన్నాయి.
ఉదాహరణకు $a\nsless a$; కనుక $\Sat{K}{\lexists[x][x<x]}$ మరియు
$\Sat/{K}{\lforall[x][\lnot x<x]}$. దీనివల్ల
$\Th{Q}\Proves/\lforall[x][\lnot x<x]$.
\end{ex}

\begin{prob}
\olref[mod][mar][mdq]{ex:model-K-of-Q}లోని $\Struct{K}$, $\Th{Q}$
యొక్క మిగిలిన స్వీకృతాలను సంతృప్తిపరుస్తుందని నిరూపించండి:
\begin{align*}
  & \lforall[x][\eq[(x \times \Obj 0)][\Obj 0]] \tag{$!Q_6$}\\
  & \lforall[x][\lforall[y][\eq[(x \times y')][((x \times y) + x)]]]
    \tag{$!Q_7$}\\
  & \lforall[x][\lforall[y][(x < y \liff
    \lexists[z][\eq[(z' + x)][y]])]] \tag{$!Q_8$}
\end{align*}
$\prime$ను మాత్రమే కలిగి, $\Struct{N}$లో సత్యమై $\Struct{K}$లో
అసత్యమయ్యే \tetoken{ఒక వాక్యాన్ని}{!!a{sentence}} కనుగొనండి.
\end{prob}

\begin{ex}
\ollabel{ex:model-L-of-Q} $\Domain{L}=\Nat\cup\{a,b\}$ గల
\tetoken{నిర్మాణం}{!!{structure}}~$\Struct{L}$ను పరిశీలించండి.
$\Assign{\prime}{L}=\nssucc$, $\Assign{+}{L}=\nsplus$ అర్థనిర్దేశాలు
ఇలా ఉన్నాయి:
\[
\begin{array}{c|c}
  x & x^\nssucc \\
  \hline
  n & n+1 \\
  a & a\\
  b & b
\end{array}
\qquad
\begin{array}{c|ccc}
  x \nsplus y & m & a & b\\
  \hline
  n & n+m & b & a\\
  a & a & b & a\\
  b & b & b & a
\end{array}
\]
$\nssucc$ \tetoken{ఒకటి-ఒకటి}{!!{injective}}, $0$ దాని పరిధిలో
లేదు, $0$ తప్ప ప్రతి $x\in\Domain{L}$ దాని పరిధిలో ఉంది కాబట్టి
$!Q_1$--$!Q_3$లు $\Struct{L}$లో సత్యం. ప్రతి $x$కూ
$x\nsplus0=x$ కాబట్టి $!Q_4$ కూడా సత్యం. $!Q_5$ కోసం
$x\nsplus y^\nssucc$, $(x\nsplus y)^\nssucc$లను పరిశీలించండి.
$x$, $y$ రెండూ ప్రామాణికమైతే అవి సమానం; ఎందుకంటే అప్పుడు
$\nssucc$, $\nsplus$లు $\prime$, $+$లతో ఏకీభవిస్తాయి. $x$
అప్రామాణికం, $y$ ప్రామాణికం అయితే
$x\nsplus y^\nssucc=x=x^\nssucc=(x\nsplus y)^\nssucc$.
$x$, $y$ రెండూ అప్రామాణికమైతే నాలుగు సందర్భాలు:
\begin{align*}
&  a \nsplus a^\nssucc  = b = b^\nssucc = (a \nsplus a)^\nssucc\\
&  b \nsplus b^\nssucc  = a = a^\nssucc = (b \nsplus b)^\nssucc\\
&  b \nsplus a^\nssucc  = b = b^\nssucc = (b \nsplus a)^\nssucc\\
&  a \nsplus b^\nssucc  = a = a^\nssucc = (a \nsplus b)^\nssucc\\
\intertext{$x$ ప్రామాణికం, $y$ అప్రామాణికం అయితే}
&  n \nsplus a^\nssucc  = n \nsplus a = b = b^\nssucc
  = (n \nsplus a)^\nssucc\\
&  n \nsplus b^\nssucc  = n \nsplus b = a = a^\nssucc
  = (n \nsplus b)^\nssucc
\end{align*}
\sourcecorrection{OLTEMODARI-009}{$b\nsplus a^\nssucc$ సందర్భంలోని
ఆంగ్ల మూల సమానతా శ్రేణి చివర $(b\nsplus y)^\nssucc$ అని, ఆ శాఖలో
నిర్దిష్టమైన $a$ స్థానంలో స్వేచ్ఛా $y$ను ఉంచుతుంది. అదే పట్టికకూ
సమానతా శ్రేణికీ సరిపడేలా $(b\nsplus a)^\nssucc$గా సరిచేశాం.}
కాబట్టి $\Sat{L}{!Q_5}$. అయితే
$a\nsplus0\neq0\nsplus a$ కాబట్టి
$\Sat/{L}{\lforall[x][\lforall[y][\eq[(x+y)][(y+x)]]]}$.
\end{ex}

\begin{prob}
\olref[mod][mar][mdq]{ex:model-L-of-Q}లోని $\Struct{L}$ను,
$\times$, $<$లకు అర్థనిర్దేశాలైన $\nstimes$, $\nsless$లను
చేర్చి విస్తరించండి. మీ నిర్మాణం $\Th{Q}$ యొక్క మిగిలిన
స్వీకృతాలను సంతృప్తిపరుస్తుందని చూపండి:
\begin{align*}
& \lforall[x][\eq[(x \times \Obj 0)][\Obj 0]] \tag{$!Q_6$}\\ &
  \lforall[x][\lforall[y][\eq[(x \times y')][((x \times y) + x)]]]
  \tag{$!Q_7$}\\ &
  \lforall[x][\lforall[y][(x < y \liff
  \lexists[z][\eq[(z'+x)][y]])]] \tag{$!Q_8$}
\end{align*}
\end{prob}

\begin{prob}
\olref[mod][mar][mdq]{ex:model-L-of-Q}లోని $\Struct{L}$లో
$a^\nssucc=a$, $b^\nssucc=b$. $a^\nssucc=b$,
$b^\nssucc=a$ అయ్యే $\Th{Q}$ నమూనా ఏదైనా ఉందా?
\end{prob}

\begin{explain}
అప్రామాణిక \tetoken{మూలకాలు}{!!{element}s}, ప్రామాణిక మూలకాలకు
“ఆవల” ఉండే $\Th{Q}$ నమూనాలను మనం స్పష్టంగా నిర్మించాం. వాస్తవానికి,
స్వీకృతాలు కనీసం అంతవరకు తప్పనిసరి చేస్తాయి. ఒక అప్రామాణిక
\tetoken{మూలకం}{!!{element}}~$x$కు ${}\nsless 0$ సంబంధం
ఉండదు; ఎందుకంటే $\Th{Q}\Proves\lforall[x][\lnot x<0]$
(\olref[inc][req][min]{lem:less-zero} చూడండి). అలాగే ప్రతి $n$కూ
$\Th{Q}\Proves\lforall[x][(x<\num{n}'\lif
(\eq[x][\num{0}]\lor\eq[x][\num{1}]\lor\dots\lor
\eq[x][\num{n}]))]$ (\olref[inc][req][min]{lem:less-nsucc}); కాబట్టి
ఏ $n>0$కూ $a\nsless n$ కాలేదు.
\end{explain}

\end{document}
